\subsection{Finite normalization of an associative family}
\label{subsec:characteristic-two-construction}

We conclude Part~\ref{part:mesh-deformations} with a general deformation-theoretic
criterion for realizing a positive filtered deformation as a specialized
fibre of an associative polynomial family (compare Gerstenhaber
\cite{Gerstenhaber}). For an algebra family over $k[t_1,\ldots,t_n]$, we call
its specialization at a parameter vector $\lambda\in k^n$ a \emph{scalar
fibre}, and we call it a \emph{binary fibre} when every coordinate of
$\lambda$ belongs to $\{0,1\}$.

\begin{theorem}
\label{thm:finite-family-normalization}
Let $\Lambda=\bigoplus_{d=0}^N\Lambda_d$ be a finite-dimensional nonnegatively
graded algebra over a field $k$ whose degree-zero part $S=\Lambda_0$ is a
finite product of copies of $k$. Let $\mu_1,\ldots,\mu_n\in C^2_S(\Lambda,\Lambda)$
be normalized $S$-relative Hochschild $2$-cocycles of positive internal
degrees $q_1,\ldots,q_n>0$. Assume that their cohomology classes span
$\HH^2(\Lambda)_{>0}$ and that the deformed product
\[
                         m_t=m+\sum_{i=1}^n t_i\mu_i
\]
is associative over the polynomial ring $k[t_1,\ldots,t_n]$, where $m$ is
the multiplication of $\Lambda$. Then every positive filtered deformation of
$\Lambda$ is isomorphic to a scalar fibre $(\Lambda,m_\lambda)$ via a
filtered isomorphism that induces the identity on $\gr\Lambda=\Lambda$ and
fixes the lifted vertex idempotents.
\end{theorem}
\begin{proof}
By Lemma~\ref{lem:split-positive-deformation}, after lifting the vertex
idempotents and splitting the filtration as $S$-bimodules, the given positive
filtered deformation is defined on the underlying graded vector space
$\Lambda$ by an associative $S$-relative multiplication $\star$ such that
$\star-m$ raises internal degree by at least $1$. We construct a parameter
vector $\lambda\in k^n$ and a sequence of filtered automorphisms of $\Lambda$
by induction on the degree $q\geq1$, starting with $q=1$ and $\lambda=0$.

Suppose that $\star$ agrees with $m_\lambda$ in all internal degrees less
than $q$, where only those coordinates $\lambda_i$ with $q_i<q$ have been
assigned. Let $\varphi_q\in C^2_S(\Lambda,\Lambda)_q$ denote the homogeneous
component of $\star-m_\lambda$ of internal degree $q$. Subtracting the
associativity identity for $m_\lambda$ from that for $\star$ in internal
degree $q$, all cross-terms involving both $\varphi_q$ and lower-degree
positive components of $\star$ or $m_\lambda$ have internal degree strictly
greater than $q$ and therefore vanish, leaving the Hochschild cocycle
condition $d\varphi_q=0$. Because the classes $[\mu_i]$ span
$\HH^2(\Lambda)_{>0}$, there exist scalars $c_i\in k$ (for those indices $i$
with $q_i=q$) and a normalized $S$-relative $1$-cochain
$f_q\in C^1_S(\Lambda,\Lambda)_q$ such that
\[
                 \varphi_q=\sum_{q_i=q}c_i\mu_i+df_q.
\]
Setting $\lambda_i=c_i$ for all $i$ with $q_i=q$ keeps $m_\lambda$
associative by hypothesis, and pulling the multiplication $\star$ back along
the filtered $S$-bimodule automorphism $1-f_q$ subtracts
$(df_q)(a,b)=a f_q(b)-f_q(ab)+f_q(a)b$ from its degree-$q$ component while
leaving all components of degree less than $q$ unchanged. Thus the
transformed product $\star$ agrees with $m_\lambda$ in all internal degrees
at most $q$. Since $\Lambda$ is finite-dimensional, $C^2_S(\Lambda,\Lambda)_q=0$
for all $q>N$, so this process terminates after at most $N$ steps and yields
the desired filtered isomorphism with $(\Lambda,m_\lambda)$.
\end{proof}

Because the residual characters $\chi_i$ extend to a $\Z$-basis of the
weight lattice $L=W(P)$ by Proposition~\ref{prop:primitive-parameter-characters},
any nonzero coordinates of the parameter vector $\lambda$ can then be rescaled
to $1$ by a diagonal grading automorphism of $\Lambda$.

\begin{corollary}
\label{cor:binary-normalization}
In the setting of Theorem~\ref{thm:finite-family-normalization}, suppose
that $\Lambda$ also carries a compatible grading by a finite-rank free
abelian group $L$ with $S$ in degree zero, and that each cocycle $\mu_i$ is
homogeneous of $L$-degree $\chi_i\in L$. If $\{\chi_1,\ldots,\chi_n\}$
extends to a $\Z$-basis of $L$, then every positive filtered deformation of
$\Lambda$ is isomorphic to a binary fibre $(\Lambda,m_\lambda)$ with
$\lambda\in\{0,1\}^n$. More generally, any two parameter vectors in $k^n$
having the same zero coordinates define isomorphic fibres.
\end{corollary}
\begin{proof}
Every group homomorphism $a:L\to k^*$ defines a grading automorphism of
$(\Lambda,m)$ that multiplies elements of $L$-weight $\ell$ by $a(\ell)$.
Transporting the deformed product $m_\lambda$ along this automorphism
rescales the $i$th parameter according to
\begin{equation}
                       \lambda'_i=\lambda_i a(\chi_i).
 \label{eq:diagonal-algebra-rescaling}
\end{equation}
Because $\{\chi_1,\ldots,\chi_n\}$ extends to a $\Z$-basis of the free
abelian group $L$, we can prescribe $a(\chi_i)\in k^*$ arbitrarily for
$i=1,\ldots,n$; in particular, setting $a(\chi_i)=\lambda_i^{-1}$ whenever
$\lambda_i\ne0$ and $a(\chi_i)=1$ when $\lambda_i=0$ transforms $m_\lambda$
into a binary fibre.
\end{proof}

Specializing Theorem~\ref{thm:finite-family-normalization} and
Corollary~\ref{cor:binary-normalization} to $A_0=kP$ yields our reduction
criterion for Theorem~A in characteristic two.

\begin{proposition}
\label{prop:family-criterion}
Let $A$ be a basic representation-finite algebra over an algebraically closed
field $k$ of characteristic two, let $A_0=kP$ ($P=\ray A_0$), and let
$I_0\subseteq I$ index the nonzero spaces $\HH^2(A_0)_{\chi_i}$ in
Proposition~\ref{prop:hochschild-mesh-restriction}. Suppose that we have
constructed normalized $S$-relative $2$-cochains
$\mu_i\in C^2_S(A_0,A_0)$ ($i\in I_0$) satisfying:
\begin{enumerate}[label={\upshape(\arabic*)}]
\item each $\mu_i$ is homogeneous of $L$-weight $\chi_i$, and its cohomology
class in $\HH^2(A_0)_{\chi_i}$ is nonzero;
\item the deformed product $m_t=m+\sum_{i\in I_0}t_i\mu_i$ is associative
over $k[t_i:i\in I_0]$;
\item in every binary fibre $(A_0,m_\lambda)$ ($\lambda\in\{0,1\}^{I_0}$),
the basis of nonzero rays $P\setminus\{0\}$ is a normed multiplicative basis,
and in every scalar fibre $(\lambda\in k^{I_0})$ the map sending each
$r\in P\setminus\{0\}$ to its ray in $(A_0,m_\lambda)$ identifies the ray
category of $(A_0,m_\lambda)$ with $P$.
\end{enumerate}
Then $A$ is isomorphic to a binary fibre $(A_0,m_\lambda)$ via a filtered
isomorphism that transports $P\setminus\{0\}$ to a normed multiplicative ray
basis of $A$ and identifies $\ray A$ with $P$ via initial forms.
\end{proposition}
\begin{proof}
As noted after Definition~\ref{def:positive-deformation}, $A$ is a positive
filtered deformation of $A_0$. By Proposition~\ref{prop:hochschild-mesh-restriction},
the cohomology classes $[\mu_i]$ ($i\in I_0$) in (1) form a basis of
$\HH^2(A_0)_{>0}$. Theorem~\ref{thm:finite-family-normalization} therefore
provides a filtered isomorphism between $A$ and a scalar fibre
$(A_0,m_{\lambda'})$ inducing the identity on $\gr A=A_0$, and
Proposition~\ref{prop:primitive-parameter-characters} together with
Corollary~\ref{cor:binary-normalization} rescales $(A_0,m_{\lambda'})$ via a
diagonal grading automorphism on the ray basis $P\setminus\{0\}$ to a binary
fibre $(A_0,m_\lambda)$. Because this diagonal rescaling preserves each
homogeneous ray line in $A_0=\gr A$, the composite filtered isomorphism
$(A_0,m_\lambda)\xrightarrow{\sim}A$ sends the normed multiplicative basis
$P\setminus\{0\}$ of (3) to a normed multiplicative basis of $A$ consisting
of one representative from each nonzero ray, and identifies $\ray A$ with
$\ray(A_0,m_\lambda)=P$ via the initial-form bijection of
Lemma~\ref{lem:ray-degree}.
\end{proof}
